\documentclass[a4paper,11pt]{article}

\usepackage[a4paper,margin=2cm]{geometry}
\usepackage{fontspec}
\usepackage{polyglossia}
\setmainlanguage{russian}
\setmainfont{Arial}
\setmonofont{Arial}

\usepackage{amsmath}
\usepackage{booktabs}
\usepackage{array}
\usepackage{tabularx}
\usepackage{enumitem}
\usepackage[table]{xcolor}
\usepackage{graphicx}
\usepackage{caption}
\usepackage{float}
\usepackage{titlesec}
\usepackage{hyperref}
\usepackage{fancyhdr}
\usepackage{tcolorbox}
\usepackage{microtype}

\definecolor{navy}{HTML}{17365D}
\definecolor{accent}{HTML}{167D9A}
\definecolor{lightblue}{HTML}{EAF2F8}
\definecolor{lightgray}{HTML}{F5F7F9}
\definecolor{warm}{HTML}{FFF6D8}
\definecolor{placeholder}{HTML}{68778A}

\hypersetup{
  colorlinks=true,
  linkcolor=navy,
  urlcolor=accent,
  citecolor=navy,
  pdftitle={LLM Classification Fine-tuning: Technical Report},
  pdfauthor={Anonymous team}
}

\captionsetup{font=small,labelfont={bf,color=navy},labelsep=period}

\titleformat{\section}{\Large\bfseries\color{navy}}{\thesection}{0.65em}{}
\titleformat{\subsection}{\large\bfseries\color{navy}}{\thesubsection}{0.55em}{}
\titlespacing*{\section}{0pt}{1.35em}{0.55em}
\titlespacing*{\subsection}{0pt}{0.95em}{0.35em}

\pagestyle{fancy}
\fancyhf{}
\setlength{\headheight}{22pt}
\lhead{\small PMLDL 2026 \textbar\ LLM Classification Fine-tuning}
\rhead{\small Technical Report}
\cfoot{\small \thepage}
\renewcommand{\headrulewidth}{0.4pt}
\renewcommand{\headrule}{\hbox to\headwidth{\color{navy}\leaders\hrule height \headrulewidth\hfill}}

\newcommand{\code}[1]{\path{#1}}
\newcommand{\metric}{\textbf{Log loss}}
\renewcommand{\arraystretch}{1.32}
\setlength{\tabcolsep}{7pt}
\arrayrulecolor{navy!45}

\newtcolorbox{importantbox}{
  colback=warm,colframe=accent!55!black,boxrule=0.55pt,arc=1.5mm,
  left=1.2mm,right=1.2mm,top=1.2mm,bottom=1.2mm
}
\newtcolorbox{visualbox}[1]{
  colback=lightgray,colframe=accent!45!white,boxrule=0.65pt,arc=1.5mm,
  height=#1,valign=center,left=4mm,right=4mm
}
\newcommand{\visualplaceholder}[2]{%
  \begin{visualbox}{#1}
    \centering
    {\large\bfseries\color{accent} Место для визуализации}\\[0.45em]
    \textcolor{placeholder}{#2}
  \end{visualbox}%
}

\begin{document}

\begin{titlepage}
\thispagestyle{empty}
\begin{center}
  \vspace*{2.7cm}
  {\Large\color{accent}\MakeUppercase{PMLDL 2026}}\\[0.8em]
  {\color{navy}\rule{0.78\textwidth}{1.1pt}}\par\vspace{1.2em}
  {\Huge\bfseries\color{navy} LLM Classification Fine-tuning}\\[0.65em]
  {\Large Технический отчёт}\\[1.25em]
  {\color{navy}\rule{0.78\textwidth}{1.1pt}}\par\vspace{2.0em}
  \begin{tcolorbox}[colback=lightblue,colframe=lightblue,width=0.84\textwidth,arc=1.5mm]
    \centering\small
    Анонимный драфт на момент фиксации финального ансамбля.\newline
    Fold 9 и публичный Kaggle score будут добавлены только после единственного
    финального запуска без изменения конфигурации.
  \end{tcolorbox}
  \vfill
  {\small \today}
\end{center}
\end{titlepage}

\section{Постановка задачи}

Для каждого примера известны пользовательский запрос и два ответа языковых
моделей. Требуется предсказать, какой вариант выберет человек: ответ Модели A,
ответ Модели B или ничью. Данные взяты из соревнования Kaggle
\cite{kaggle}. Модель возвращает три вероятности, соответствующие этим
вариантам.

Основная метрика --- multiclass \metric:

\begin{equation}
  \mathcal{L} = -\frac{1}{N}\sum_{i=1}^{N}\sum_{k=1}^{3}
  y_{ik}\log p_{ik},
  \label{eq:logloss}
\end{equation}

где $y_{ik}$ равно единице, если для $i$-го примера верен класс $k$, и нулю
иначе, а $p_{ik}$ --- вероятность этого класса. Меньшее значение означает
лучшее качество и более точную калибровку вероятностей.

\begin{table}[H]
\centering
\caption{Фиксированный протокол разбиения данных.}
\label{tab:protocol}
\small
\begin{tabularx}{\textwidth}{|>{\bfseries}p{0.23\textwidth}|p{0.18\textwidth}|>{\raggedleft\arraybackslash}p{0.16\textwidth}|X|}
\hline
\rowcolor{lightblue} Часть данных & Folds & Примеров & Назначение \\ \hline
Обучение & 0--6 & 40\,235 & Обучение моделей. \\ \hline
Выбор кандидатов & 7 & 5\,746 & Сравнение baseline-моделей, абляции и выбор состава ансамбля. \\ \hline
Калибровка & 8 & 5\,748 & Подбор температуры и весов уже выбранных моделей. \\ \hline
Финальная оценка & 9 & 5\,748 & Однократная оценка только после фиксации конфигурации. \\ \hline
\end{tabularx}
\end{table}

\begin{figure}[H]
\centering
\visualplaceholder{3.3cm}{
Схема задачи: запрос, ответ Модели A и ответ Модели B поступают в классификатор,\\
который возвращает вероятности выбора A, B и ничьей.}
\caption{Пайплайн постановки задачи.}
\label{fig:task-pipeline}
\end{figure}

\section{Baseline-модели}

Все модели в таблице обучались на folds 0--6 и сравнивались на fold 7. Поэтому
различия в \metric{} отражают различия в моделях, а не в выборках. Простая
логистическая регрессия использует структурные признаки диалога. Второй
baseline добавляет TF-IDF признаки слов и символьных $n$-грамм. Улучшенный
sparse TF-IDF получен последовательным изменением одного параметра за раз.

\begin{table}[H]
\centering
\caption{Baseline-результаты на fold 7.}
\label{tab:baselines}
\small
\begin{tabularx}{\textwidth}{|X|>{\raggedleft\arraybackslash}p{0.16\textwidth}|>{\raggedleft\arraybackslash}p{0.16\textwidth}|>{\raggedleft\arraybackslash}p{0.16\textwidth}|}
\hline
\rowcolor{lightblue} Модель & \metric & Accuracy & Macro F1 \\ \hline
Structural Logistic Regression & 1.06303 & 0.44692 & 0.44014 \\ \hline
Word + character TF-IDF & 1.04052 & 0.46920 & 0.46153 \\ \hline
Улучшенный sparse TF-IDF & 1.03576 & 0.47355 & 0.46595 \\ \hline
\end{tabularx}
\end{table}

Лучший sparse baseline снизил log loss на $0.00476$ по сравнению с исходным
word + character TF-IDF. Это подтвердило, что разреженная модель полезна как
самостоятельный кандидат и как компонент будущего ансамбля.

\begin{figure}[H]
\centering
\visualplaceholder{4.0cm}{
Горизонтальная диаграмма log loss для трёх baseline-моделей.\\
Ось $x$: log loss, меньше --- лучше.}
\caption{Сравнение baseline-моделей.}
\label{fig:baseline-comparison}
\end{figure}

\section{Предлагаемая нейросетевая модель}

Нейросетевой кандидат --- \texttt{microsoft/deberta-v3-base}, дообученный
QLoRA. Вход формируется из prompt, ответа Модели A и ответа Модели B. Каждый
текст сериализуется по ходам диалога; затем применяется balanced head-tail
truncation. Во время обучения ответы A и B случайно меняются местами с
вероятностью $0.5$, а во время inference вероятности для исходного и
переставленного примеров усредняются. Это уменьшает чувствительность модели к
порядку ответов.

\begin{table}[H]
\centering
\caption{Конфигурация DeBERTa QLoRA, seed 42.}
\label{tab:deberta-card}
\small
\begin{tabularx}{\textwidth}{|>{\bfseries}p{0.34\textwidth}|X|}
\hline
\rowcolor{lightblue} Параметр & Значение \\ \hline
Backbone & \texttt{microsoft/deberta-v3-base}, revision \texttt{8ccc9b6f...}. \\ \hline
Способ обучения & QLoRA в 4-bit NF4; double quantization. \\ \hline
Адаптеры & $r=16$, $\alpha=32$, dropout $=0.05$; \texttt{query\_proj} и \texttt{value\_proj}. \\ \hline
Обучаемые модули & Адаптеры, \texttt{pooler} и \texttt{classifier}. \\ \hline
Длина входа & Максимум 512 токенов. \\ \hline
Оптимизация & AdamW, learning rate $2.8\cdot10^{-4}$, weight decay $0.01$, 3 эпохи. \\ \hline
Batch и scheduler & Effective batch size 32; warmup 0.06; scheduler horizon 5 эпох. \\ \hline
Симметрия A/B & Swap augmentation при обучении и усреднение вероятностей при inference. \\ \hline
\end{tabularx}
\end{table}

На fold 7 seed 42 получил \metric{} $1.01505$, accuracy $0.48747$ и macro F1
$0.48420$. Это лучший отдельный нейросетевой кандидат из проверенных seed.

\begin{figure}[H]
\centering
\visualplaceholder{5.1cm}{
Архитектурная схема: сериализация диалога $\rightarrow$ tokenizer и truncation\\
$\rightarrow$ DeBERTa encoder с QLoRA adapters $\rightarrow$ classification head\\
$\rightarrow$ вероятности выбора A, B и ничьей.}
\caption{Архитектура нейросетевого кандидата.}
\label{fig:deberta-architecture}
\end{figure}

\section{Улучшения и контролируемые эксперименты}

Все изменения оценивались на том же fold 7. Для sparse модели проводился
последовательный one-factor-at-a-time sweep: после каждого шага менялся ровно
один параметр TF-IDF или классификатора, а победитель становился контролем для
следующего шага. Для DeBERTa зафиксированная конфигурация QLoRA запускалась с
тремя seed; затем сравнивались отдельные seed и их среднее вероятностей.

\begin{table}[H]
\centering
\caption{Подтверждённые улучшения на fold 7.}
\label{tab:improvements}
\small
\begin{tabularx}{\textwidth}{|p{0.27\textwidth}|X|>{\raggedleft\arraybackslash}p{0.16\textwidth}|>{\raggedleft\arraybackslash}p{0.16\textwidth}|}
\hline
\rowcolor{lightblue} Эксперимент & Изменение & \metric & Изменение log loss \\ \hline
Sparse TF-IDF sweep & Последовательный подбор $n$-грамм, словаря, \texttt{min\_df} и $\alpha$. & 1.03576 & $-0.00476$ относительно B2. \\ \hline
DeBERTa QLoRA, seed 42 & Выбор лучшего seed при неизменной QLoRA-конфигурации. & 1.01505 & Лучший отдельный seed. \\ \hline
Среднее трёх seed & Усреднение seed 17, 42 и 73. & 1.01889 & $+0.00384$ относительно seed 42. \\ \hline
\end{tabularx}
\end{table}

\begin{table}[H]
\centering
\caption{Влияние seed на DeBERTa QLoRA, fold 7.}
\label{tab:seeds}
\small
\begin{tabularx}{\textwidth}{|>{\bfseries}p{0.22\textwidth}|>{\raggedleft\arraybackslash}p{0.18\textwidth}|>{\raggedleft\arraybackslash}p{0.18\textwidth}|>{\raggedleft\arraybackslash}p{0.18\textwidth}|X|}
\hline
\rowcolor{lightblue} Конфигурация & \metric & Accuracy & Macro F1 & Решение \\ \hline
Seed 17 & 1.02800 & 0.46711 & 0.46777 & Не выбран. \\ \hline
Seed 42 & 1.01505 & 0.48747 & 0.48420 & Выбран для ансамбля. \\ \hline
Seed 73 & 1.02807 & 0.47355 & 0.47328 & Не выбран. \\ \hline
Среднее трёх seed & 1.01889 & 0.48016 & 0.48001 & Уступает seed 42. \\ \hline
\end{tabularx}
\end{table}

\begin{figure}[H]
\centering
\visualplaceholder{3.9cm}{
Dot plot или bar chart: log loss для seed 17, 42, 73 и их среднего.\\
Выделить выбранный seed 42.}
\caption{Стабильность DeBERTa QLoRA по seed.}
\label{fig:seed-stability}
\end{figure}

\section{Финальный ансамбль и результаты}

Состав ансамбля был выбран на fold 7: DeBERTa QLoRA seed 42 и улучшенный sparse
TF-IDF. На fold 8 проверены три набора положительных весов и для каждого
подобрана одна температура. Победила равная смесь вероятностей с последующим
temperature scaling, $T=0.817311$. Ни Kaggle score, ни fold 9 не использовались
при выборе моделей, весов или температуры.

\begin{table}[H]
\centering
\caption{Сравнение на calibration fold 8.}
\label{tab:calibration-comparison}
\small
\begin{tabularx}{\textwidth}{|p{0.30\textwidth}|>{\raggedleft\arraybackslash}p{0.12\textwidth}|>{\raggedleft\arraybackslash}p{0.11\textwidth}|>{\raggedleft\arraybackslash}p{0.11\textwidth}|>{\raggedleft\arraybackslash}p{0.11\textwidth}|}
\hline
\rowcolor{lightblue} Модель & \metric & Accuracy & Macro F1 & ECE \\ \hline
Калиброванная DeBERTa QLoRA & 1.00892 & 0.48713 & 0.48349 & 0.01833 \\ \hline
Калиброванный sparse TF-IDF & 1.02609 & 0.49791 & 0.49214 & 0.03329 \\ \hline
Финальный ансамбль & \textbf{0.99697} & \textbf{0.50035} & \textbf{0.49521} & \textbf{0.01111} \\ \hline
\end{tabularx}
\end{table}

\begin{table}[H]
\centering
\caption{Проверенные веса ансамбля после temperature scaling, fold 8.}
\label{tab:weight-grid}
\small
\begin{tabularx}{\textwidth}{|X|>{\raggedleft\arraybackslash}p{0.20\textwidth}|>{\raggedleft\arraybackslash}p{0.16\textwidth}|>{\raggedleft\arraybackslash}p{0.16\textwidth}|}
\hline
\rowcolor{lightblue} Веса (DeBERTa, sparse) & Температура & \metric & ECE \\ \hline
0.25, 0.75 & 0.87399 & 1.00397 & 0.01542 \\ \hline
0.50, 0.50 & \textbf{0.81731} & \textbf{0.99697} & \textbf{0.01111} \\ \hline
0.75, 0.25 & 0.84529 & 0.99899 & 0.01746 \\ \hline
\end{tabularx}
\end{table}

\begin{importantbox}
\textbf{Интерпретация результатов.} Значение 0.99697 получено на fold 8,
использованном для калибровки и выбора весов. Его нельзя называть финальным
обобщающим результатом. Fold 9 остаётся закрытым до запуска
\code{scripts/finalize_release.py} с уже зафиксированным
\code{configs/final_model.json}.
\end{importantbox}

\begin{figure}[H]
\centering
\visualplaceholder{4.2cm}{
Reliability diagram до и после temperature scaling для финального ансамбля.\\
В подписи указать ECE до калибровки 0.03190 и после калибровки 0.01111.}
\caption{Калибровка вероятностей ансамбля на fold 8.}
\label{fig:calibration}
\end{figure}

\begin{figure}[H]
\centering
\visualplaceholder{4.8cm}{
Схема: вероятности DeBERTa и sparse TF-IDF $\rightarrow$ веса 0.5 и 0.5\\
$\rightarrow$ temperature scaling $T=0.817311$ $\rightarrow$ итоговые вероятности.}
\caption{Архитектура финального ансамбля.}
\label{fig:ensemble-architecture}
\end{figure}

\subsection{Итоговая оценка и Kaggle}

На момент подготовки этого драфта fold 9 не открывался, поэтому итоговый log
loss и Kaggle public score здесь намеренно не указаны. После получения
вероятностей двух зафиксированных моделей на fold 9 и test-наборе команда
запустит \code{make finalize}; команда создаст метрики fold 9 и
\code{results/final/submission.csv}. В отчёт будут добавлены только значения
из \code{results/final/final_holdout_metrics.json} и подтверждённый public
score отправленного файла.

\section{Воспроизводимость}

Репозиторий фиксирует разбиение, конфигурации моделей, результаты выбора и
конфигурацию финального ансамбля. Большие файлы не включены в Git: для них
указаны контрольные суммы и внешние ссылки. Это отделяет воспроизводимый код
от лицензированных данных и чекпойнтов.

\begin{table}[H]
\centering
\caption{Артефакты воспроизводимого release.}
\label{tab:reproducibility}
\small
\begin{tabularx}{\textwidth}{|>{\bfseries}p{0.34\textwidth}|X|}
\hline
\rowcolor{lightblue} Артефакт & Назначение \\ \hline
\code{configs/split.json} и \code{data/splits/folds.csv} & Фиксированные роли folds и назначение примеров. \\ \hline
\code{configs/experiments/} & Конфигурации retained baseline, sparse и DeBERTa кандидатов. \\ \hline
\code{configs/final_model.json} & Замороженные модели, веса $0.5/0.5$ и температура $0.817311$. \\ \hline
\code{results/model_comparison.csv} & Компактная таблица результатов, используемых в отчёте. \\ \hline
\code{results/final/} & Evidence выбора кандидатов, перебора весов и калибровки. \\ \hline
\code{artifacts/FINAL_ARTIFACTS.md} & SHA-256 и расположение двух внешних чекпойнтов. \\ \hline
\code{docs/REPRODUCE.md} & Порядок восстановления артефактов и запуска финального pipeline. \\ \hline
\code{requirements*.lock} & Зафиксированные зависимости. \\ \hline
\end{tabularx}
\end{table}

\begin{figure}[H]
\centering
\visualplaceholder{4.5cm}{
Схема воспроизведения: data audit $\rightarrow$ frozen folds $\rightarrow$ train\\
$\rightarrow$ selection на fold 7 $\rightarrow$ calibration на fold 8\\
$\rightarrow$ один запуск fold 9 $\rightarrow$ \texttt{submission.csv}.}
\caption{Воспроизводимый pipeline.}
\label{fig:reproducibility-pipeline}
\end{figure}

\section{Заключение}

На фиксированном selection fold лучшим отдельным кандидатом стала DeBERTa QLoRA
seed 42 с log loss 1.01505. Улучшенный sparse TF-IDF оказался достаточно
дополняющим её кандидатом для ансамбля. На calibration fold равный ансамбль с
temperature scaling уменьшил log loss до 0.99697 и ECE до 0.01111.

Окончательное утверждение о качестве будет сделано после единственной оценки
на fold 9. Основные ограничения текущей оценки --- зависимость выбора от одного
selection fold, вариативность между seed и отсутствие ещё не открытого
независимого holdout-результата.

{\footnotesize
\clearpage
\begin{thebibliography}{9}
\bibitem{kaggle}
Kaggle. \emph{LLM Classification Fine-tuning}.\\
\url{https://www.kaggle.com/competitions/llm-classification-finetuning}

\bibitem{deberta}
P. He, X. Liu, J. Gao, and W. Chen. DeBERTa: Decoding-enhanced BERT with
Disentangled Attention. ICLR, 2021.\\
\url{https://arxiv.org/abs/2006.03654}

\bibitem{lora}
E. Hu et al. LoRA: Low-Rank Adaptation of Large Language Models. ICLR, 2022.\\
\url{https://arxiv.org/abs/2106.09685}

\bibitem{qlora}
T. Dettmers et al. QLoRA: Efficient Finetuning of Quantized LLMs. NeurIPS, 2023.\\
\url{https://arxiv.org/abs/2305.14314}
\end{thebibliography}
}

\end{document}
